\documentclass{article}
\usepackage[margin=1in]{geometry}
\usepackage{amsmath, amssymb, amsthm}
\usepackage{enumitem}

%Formatting and Spacing
\setenumerate[1]{label = (\alph*), parsep = 1pt, topsep = 5pt}
\setlength\parindent{0pt}
\linespread{1.15}

%Titles and Headers Custom Fields
\newcommand{\lectTitle}{Lecture 1 Notes}
\newcommand{\lectTime}{August 23, 2021}
\newcommand{\lectClass}{Honors Discrete Mathematics}
\newcommand{\lectClassInstructor}{Professor Gerandy Brito}
\newcommand{\lectSection}{Fall 2021}
\newcommand{\lectAuthorName}{Sarthak Mohanty}

%Fancy Headers and Footers
\usepackage{fancyhdr}
\usepackage{extramarks}
\pagestyle{fancy}
\lhead{\lectTime}
\chead{\lectClass \ (\lectClassInstructor)}
\rhead{\lectTitle}
\cfoot{\thepage}
\renewcommand\headrulewidth{0.4pt}
\renewcommand\footrulewidth{0.4pt}

\title{
    \vspace{2in}
    \textbf{\lectClass:\\ \lectTitle}\\
    % \normalsize\vspace{0.1in}\small{Due\ on\ \lectDueDate\ at 11:59pm}\\
    \vspace{0.1in}\large{\textit{\lectClassInstructor\ \lectSection}}
    \vspace{3in}
    \author{\textbf{\lectAuthorName}}
    \date{}
}

\begin{document}

\maketitle
\pagebreak

\section{Logic}

The goal of discrete mathematics is to translate our current model of information (English, as an example), into an universal language. We call this language Logic. The precise definition of ``logic" is quite broad. When most people say ``logic", they mean either propositional logic or predicate logic. Let us begin by introducing propositional logic.

\subsection{Propositions}
    To begin, we have the following definition.
    
    \vspace{1.5mm}
    \textbf{Definition 1.1.1} \textit{A proposition is a statement that can be always either true or false  (never both).}
    
    \vspace{1.5mm}
    A few examples:
    \begin{itemize}[itemsep = 0.2mm, parsep = 5pt, topsep = 5pt]
        \item Professor Brito is Cuban. (True)
        \item 2 + 3 = GT. (False)
        \item ``If we go to the park, then we play baseball." Both the antecedent and the consequent are propositions, so the entire statement is True.
    \end{itemize}
    
    An example of a statement that is \textit{not} a proposition  is the statement ``CS 2051 is a fun class!", because there is no objective definition for what `fun' is. Let's conclude with one more definition.
    
    \vspace{1.5mm}
    \textbf{Definition 1.1.2}: \textit{A theorem is a proposition that is always true but not immediate to verify.}

\subsection{Logical Operators}
    First, some notation:
    \begin{itemize}[itemsep = 0.2mm, parsep = 5pt, topsep = 5pt]
        \item Propositions will generally be labeled as $p, q, r, \dots$.
        \item The negation of a proposition is indicated by the $\neg$ symbol. (ex. $\neg p$).
        \item The conjunction of two propositions is indicated by the symbol $\land$.
        \item The disjunction of two propositions is indicated by the symbol $\lor$.
        \item The statement ``If $p$ then $q$" can be represented by the notation $p \Rightarrow q$. This is also known as a conditional statement.
    \end{itemize}
    Using this newly defined notation, let's introduce a new proposition. $$ \text{For any natural number $n$, and for any set of integers $x$, $y$, and $z$, we have $x^n + y^n = z^n \Rightarrow xyz = 0$.}$$ The truth value of this proposition is clearly false (choose $n = 1$ and corresponding $x, y, z$). However, restricting the domain of $n$ to $n > 2$, this proposition becomes true (see Fermat's Last Theorem).
    
    \vspace{1.5mm}
    Truth tables can also be used to more easily determine the validity of a propositional statement.
    $$\begin{array}{|c|c|c|c|c|c|}
        \hline
        p & q & \neg p & p \land q & p \lor q & p \Rightarrow q\\
        \hline
        T & T & F & T & F & T \\
        T & F & F & F & T & F \\
        F & T & T & F & T & T \\
        F & F & T & F & F & T \\
        \hline
    \end{array}$$

\subsection{Direct Proofs}
    The goal of a direct proof is to determine the truth value of a proposition with as few assumptions as possible.

    \vspace{1.5mm}
    \textbf{Example}
    
    A natural number $n$ is said to be \textit{even} if there is another natural number $k$ such that $n = 2k$. Show that if $n$ and $m$ are even numbers, then $n + m$ is even.
    
    \vspace{1.5mm}
    \textbf{Solution}
    
    Suppose $n$ and $m$ are even numbers. Then there exists natural numbers $k$ and $l$ such that $n = 2k$ and $m = 2l$. Now $n + m = 2(k + l)$, so by definition $n + m$ is even.

\end{document}